\documentclass[a4paper, twoside]{article}
\usepackage[letterpaper, top=0.6in, bottom=0.6in, left=1in, right=1in]{geometry}

\renewcommand{\baselinestretch}{1.15}
\raggedbottom
\setlength{\parindent}{0pt}
\setlength{\parskip}{0.8em}
\renewcommand{\thefootnote}{\roman{footnote}} 
\usepackage{titlesec}
\usepackage{enumitem}
\setlist{parsep=2pt, topsep=0pt, partopsep=1pt, itemsep=1pt, leftmargin=10mm}
\usepackage{fancyhdr}
\usepackage{lipsum}
\usepackage{etoolbox}
\usepackage{tcolorbox}
\usepackage{makecell}
\usepackage[table]{xcolor}
\usepackage{nicematrix}
\usepackage{media9}
\usepackage{array}
\usepackage{hyperref}
\hypersetup{colorlinks=true, linkcolor=blue, urlcolor=blue}
\usepackage{chngcntr}
\usepackage{multicol}
\usepackage{multirow}
\usepackage{graphicx}
\usepackage{amsfonts}
\usepackage{amsmath}
\usepackage{amssymb}
\usepackage{amsthm}
\usepackage[dvipsnames]{xcolor}
\usepackage{float}
\usepackage{pgfplots}
\usepackage{pgfplotstable}
\usetikzlibrary{3d}
\pgfplotsset{compat=1.18}
\pagestyle{fancy}
\usepackage{titlesec}
\titlespacing*{\section}
  {0pt}
  {1.5em plus 1em minus .5em}
  {0ex}
\usepackage{changepage}


\fancypagestyle{standard}{
    \fancyhf{}
    \fancyhead[R]{\footnotesize\bfseries\thepage}
}

\fancypagestyle{copyright}{
    \fancyhf{} 
    \fancyfoot[R]{\footnotesize ©~2026 Archimedes314 Notes---MIT License}
    \fancyhead[R]{\footnotesize\bfseries\thepage}
}

\renewcommand{\headrulewidth}{0pt}
\renewcommand{\footrulewidth}{0pt}

\newcommand{\copyrightpage}{
  \pagestyle{copyright}
}

\newcommand{\standardpage}{
    \pagestyle{standard}
}


\titleformat{\section}
  {\normalfont\Large\bfseries}
  {\thesection}{1em}{}

\renewcommand{\sectionmark}[1]{\markboth{#1}{}}

\newcounter{definitioncounter}
\renewcommand{\thedefinitioncounter}{\romannumeral\value{definitioncounter}}
\newcommand{\defc}{\refstepcounter{definitioncounter}\thedefinitioncounter}
\newcommand{\defcs}{\thesection.\refstepcounter{definitioncounter}\thedefinitioncounter}
\newcommand{\defcss}{\thesubsection.\refstepcounter{definitioncounter}\thedefinitioncounter}
\newcommand{\defcsss}{\thesubsubsection.\refstepcounter{definitioncounter}\thedefinitioncounter}

\newcounter{examplecounter}
\renewcommand{\theexamplecounter}{\romannumeral\value{examplecounter}}
\newcommand{\egc}{\refstepcounter{examplecounter}\theexamplecounter  }
\newcommand{\egcs}{\thesection.\refstepcounter{examplecounter}\theexamplecounter  }
\newcommand{\egcss}{\thesubsection.\refstepcounter{examplecounter}\theexamplecounter  }
\newcommand{\egcsss}{\thesubsubsection\refstepcounter{examplecounter}\theexamplecounter  }

\newcounter{theoremcounter}
\renewcommand{\thetheoremcounter}{\romannumeral\value{theoremcounter}}
\newcommand{\thc}{\refstepcounter{theoremcounter}\thetheoremcounter  }
\newcommand{\thcs}{\thesection\refstepcounter{theoremcounter}\thetheoremcounter  }
\newcommand{\thcss}{\thesubsection\refstepcounter{theoremcounter}\thetheoremcounter  }
\newcommand{\thcsss}{\thesubsubsection\refstepcounter{theoremcounter}\thetheoremcounter  }

\newcounter{examplecounterrefer}
\renewcommand{\theexamplecounter}{\romannumeral\value{examplecounter}}
\newcommand{\egr}[1]{(\ref{#1})}
\newcommand{\egrs}[1]{(\thesection.\ref{#1})}
\newcommand{\egrss}[1]{(\thesubsection.\ref{#1})}
\newcommand{\egrsss}[1]{(\thesubsubsection.\ref{#1})}

\newcommand{\mc}[2]{\multicolumn{#1}{c}{#2}}

%Roman and bold font shortcuts
\renewcommand{\Re}{\mathrm{Re}}
\renewcommand{\Im}{\mathrm{Im}}
\newcommand{\diam}{\mathrm{diam}}
\newcommand{\cis}{\mathrm{cis}}
\newcommand{\Arg}{\mathrm{Arg}}
\renewcommand{\i}{\mathrm{i}}
\renewcommand{\d}{\mathrm{d}}
\newcommand{\D}{\mathrm{D}}
\newcommand{\e}{\mathrm{e}}
\newcommand{\C}{\mathbb{C}}
\newcommand{\Ce}{\hat{\mathbb{C}}}
\newcommand{\R}{\mathbb{R}}
\newcommand{\Q}{\mathbb{Q}}
\newcommand{\Z}{\mathbb{Z}}
\newcommand{\N}{\mathbb{N}}
\renewcommand{\S}{\mathrm{S}}
\newcommand{\textdef}[1]{\hspace{0.2cm}\fontsize{11}{5}\selectfont\textbf{#1}\normalsize}


\renewcommand{\hat}{\widehat}
\renewcommand{\bar}[1]{\overline{#1}}
\newcommand{\ttilde}[1]{\underset{\Tilde{}}{#1}}
\newcommand{\vtilde}[1]{\underset{\Tilde{}}{\mathbf{#1}}}
\newcommand{\vect}[1]{\overrightharpoon{\mathbf{#1}}}
\newcommand{\rarrow}{\rightarrow}
\newcommand{\lorarrow}{\longrightarrow}
\newcommand{\ioi}{\longleftrightarrow}
\newcommand{\sioi}{\leftrightarrow}
\newcommand{\defeq}{\overset{\mathrm{def}}{=}}
\newcommand{\geqs}{\geqslant}
\newcommand{\leqs}{\leqslant}
\newcommand{\nsubset}{\not\subset}

\makeatletter
\DeclareFontFamily{U}{MnSymbolA}{}
\DeclareFontShape{U}{MnSymbolA}{m}{n}{
    <-6>  MnSymbolA5
   <6-7>  MnSymbolA6
   <7-8>  MnSymbolA7
   <8-9>  MnSymbolA8
   <9-10> MnSymbolA9
  <10-12> MnSymbolA10
  <12->   MnSymbolA12}{}
\DeclareFontShape{U}{MnSymbolA}{b}{n}{
    <-6>  MnSymbolA-Bold5
   <6-7>  MnSymbolA-Bold6
   <7-8>  MnSymbolA-Bold7
   <8-9>  MnSymbolA-Bold8
   <9-10> MnSymbolA-Bold9
  <10-12> MnSymbolA-Bold10
  <12->   MnSymbolA-Bold12}{}
\DeclareSymbolFont{MnSyA}{U}{MnSymbolA}{m}{n}
\SetSymbolFont{MnSyA}{bold}{U}{MnSymbolA}{b}{n}

\DeclareRobustCommand{\overleftharpoon}{\mathpalette{\overarrow@\leftharpoonfill@}}
\DeclareRobustCommand{\overrightharpoon}{\mathpalette{\overarrow@\rightharpoonfill@}}
\def\leftharpoonfill@{\arrowfill@\leftharpoondown\mn@relbar\mn@relbar}
\def\rightharpoonfill@{\arrowfill@\mn@relbar\mn@relbar\rightharpoonup}

\DeclareMathSymbol{\leftharpoondown}{\mathrel}{MnSyA}{'112}
\DeclareMathSymbol{\rightharpoonup}{\mathrel}{MnSyA}{'100}
\DeclareMathSymbol{\mn@relbar}{\mathrel}{MnSyA}{'320}
\makeatother

\makeatletter
\def\@afterindentfalse{\setcounter{examplecounter}{0}}
\makeatother

\makeatletter
\def\mathcolor#1#{\@mathcolor{#1}}
\def\@mathcolor#1#2#3{%
  \protect\leavevmode
  \begingroup
    \color#1{#2}#3%
  \endgroup
}
\makeatother


\begin{document}
\numberwithin{equation}{subsubsection}
\renewcommand{\arraystretch}{1.5}

\newcommand{\bluebox}[2]{\begin{tcolorbox}[colback=blue!5!white, colframe=blue!75!black, title=#1] #2 \end{tcolorbox}}

\newcommand{\egbluebox}[2]{
    \begin{tcolorbox}[
        colback=blue!5!white, 
        colframe=blue!75!black, 
        title=Example \egc\label{#1}
    ] 
    #2 \end{tcolorbox}
}

\newcommand{\egblueboxs}[2]{
    \begin{tcolorbox}[
        colback=blue!5!white, 
        colframe=blue!75!black, 
        title=Example \thesection.\egc\label{#1}
    ]   
    #2 \end{tcolorbox}
}

\newcommand{\egblueboxss}[2]{
    \begin{tcolorbox}[
        colback=blue!5!white, 
        colframe=blue!75!black, 
        title=Example 
        \thesubsection.\egc\label{#1}
    ] 
    #2 \end{tcolorbox}
}

\newcommand{\egblueboxsss}[2]{
    \begin{tcolorbox}[
        colback=blue!5!white, 
        colframe=blue!75!black, 
        title=Example 
        \thesubsubsection.\egc\label{#1}
    ] 
    #2 \end{tcolorbox}
}

\newcommand{\aInt}[2]{\int_{#1}^{#2}}
\newcommand{\aIInt}[4]{\int_{#1}^{#2} \int_{#3}^{#4}}
\newcommand{\bInt}[1]{\int_{#1}}
\newcommand{\bIInt}[2]{\int_{#1}\int_{#2}}
\newcommand{\Int}[4]{\int_{#1}^{#2}{#3\,\text{d} #4}}
\newcommand{\IInt}[7]{\int_{#1}^{#2}{\int_{#3}^{#4}{#5\,\text{d} #6\,\text{d} #7}}}
\newcommand{\lInt}[3]{\int_{#1}{#2\,\d #3}}
\newcommand{\mlInt}[3]{\int_{#1}{\left|#2\right|\,|\d #3|}}
\newcommand{\lIInt}[3]{\iint\limits_{#1}{#2\,\d #3}}
\newcommand{\lIIInt}[3]{\iiint\limits_{#1}{#2\,\d #3}}
\newcommand{\Oint}[3]{\oint_{#1}{#2\,\d #3}}
\newcommand{\mOint}[3]{\oint_{#1}{|#2|\,|\d #3|}}
\newcommand{\nfoldInt}[5]{{}_{#1}\!\int_{#2}^{#3}{#4\, \d #5^{#1}}}
\newcommand{\diff}[2]{\frac{\text{d} #1}{\text{d} #2}}
\newcommand{\pdiff}[2]{\frac{\partial #1}{\partial #2}}
\newcommand{\ndiff}[3]{\frac{\d^{#1}#2}{\d {#3}^{#1}}}
\newcommand{\EulerD}[4]{{}_{#2}^{#1}{\mathrm{D}}_{#4}^{#3}}
\newcommand{\npartialone}[3]{\frac{\partial^{#1}#2}{\partial {#3}^{#1}}}
\newcommand{\npartialtwo}[6]{\frac{\partial^{#1}{#4}}{\partial{#5}^{#2} \partial{#6}^{#3}}}
\newcommand{\npartialthree}[8]{\frac{\partial^{#1}{#5}}{\partial{#6}^{#2} \partial{#7}^{#3} \partial{#8}^{#4}}}
\newcommand{\tvec}[1]{\underset{\Tilde{}}{\mathbf{#1}}}
\newcommand{\vvec}[2]{\langle #1, #2 \rangle}
\newcommand{\vvvec}[3]{\langle #1, #2, #3 \rangle}
\newcommand{\VVec}[2]{\left[\begin{array}{c}#1\\#2\end{array}\right]}
\newcommand{\VVVec}[3]{\left[\begin{array}{c}#1\\#2\\#3\end{array}\right]}
\newcommand{\norm}[1]{\left\|#1\right\|}
\newcommand{\normone}[1]{\left\|#1\right\|_1}
\newcommand{\normtwo}[1]{\left\|#1\right\|_2}
\newcommand{\norminf}[1]{\left\|#1\right\|_\infty}
\newcommand{\modu}[1]{\left|#1\right|}
\newcommand{\brackets}[1]{\left(#1\right)}
\newcommand{\sqbrackets}[1]{\left[#1\right]}
\newcommand{\suma}[3]{\sum\limits_{#1}^{#2} #3}
\newcommand{\sumb}[2]{\sum\limits_{#1} #2}
\newcommand{\proda}[3]{\prod\limits_{#1}^{#2} #3}
\newcommand{\prodb}[2]{\prod\limits_{#1} #2}
\newcommand{\Lim}[3]{\lim_{#1\to #2} #3}
\newcommand{\Limb}[3]{\lim_{#1\to #2} \left(#3\right)}
\newcommand{\Lims}[3]{\lim_{#1\to #2} \left[#3\right]}
\newcommand{\Max}[2]{\max_{#1}{#2}}
\newcommand{\To}[2]{\overset{#1\to#2}{\lorarrow}}
\newcommand{\Res}[1]{\underset{#1}{\text{Res}}}
\newcommand{\Log}[2]{\log_{#1}{#2}}
\newcommand{\Logb}[2]{\log_{#1}{\left(#2\right)}}
\newcommand{\Logs}[2]{\log_{#1}{\left[#2\right]}}
\newcommand{\floor}[1]{\left\lfloor #1 \right\rfloor}
\newcommand{\ceil}[1]{\left\lceil #1 \right\rceil}
\newcommand{\eval}[3]{\left[ #1\right|_{#2}^{#3}}
\newcommand{\evalone}[3]{\left( #1\right|_{#2}^{#3}}
\newcommand{\evaltwo}[3]{\left. #1\right|_{#2}^{#3}}
\newcommand{\sssu}[1]{\subsubsection*{\underline{#1}}}
\pagenumbering{gobble}

\begin{titlepage}
\thispagestyle{copyright}
\centering

\vspace*{3cm}

{\Huge\scshape Syntax \& Semantics\par}
\vspace{0.3cm}

{\large Topic 1, Section 1\par}

\vspace{1.2cm}

{\LARGE\bfseries Complete Section \\[0.3em]
\large Overview\par}
\vfill

\begin{minipage}{0.7\textwidth}
\centering
\small\itshape
`Give me a lever long enough and a fulcrum on which to place it and I shall move the world.'\\[0.3cm]
— Archiemdes (c. 287--212 BC)
\end{minipage}

\vspace{1.5cm}

\end{titlepage}

\newpage

\tableofcontents

\newpage

\newpage
\standardpage

\restoregeometry 
\pagenumbering{arabic}
\setcounter{page}{1}

\section{Section Overview}
This section focuses on the syntax (structure and symbols) and semantics (meaning and truth) of propositional logic. We will see how statements are built using logical connectives and how their truth values are determined.

Propositional logic is a formal system built from atomic statements (propositions) and logical connectives, such as 'and', 'or', 'not', and 'if...then'. The syntax of propositional logic defines the rules for combining these symbols into valid formulas, while the semantics assigns truth values to each formula based on the truth of its parts. Mastering syntax means knowing how to write well-formed formulas, and mastering semantics means understanding what those formulas mean and when they are true or false.

We will cover the following key aspects:

\begin{itemize}
    \item \textbf{Syntax}: What counts as a proposition? How do we build more complex statements from simpler ones? What are the rules for writing formulas?

    \item \textbf{Semantics}: How do we assign truth values to statements? What does it mean for a formula to be true, false, or always true (a tautology)?

    \item \textbf{Logical equivalence}: When do two statements mean the same thing?

    \item \textbf{Satisfiability}: When is it possible for a formula to be true?
\end{itemize}

Understanding these principles is crucial for constructing clear mathematical arguments and for all future work in logic and mathematics.

It is also helpful to know how propositional logic compares to other types of logic. Propositional logic deals only with whole statements (propositions) that can be either true or false. In contrast, other forms of logic allow us to express more detailed relationships:

\begin{itemize}
    \item \textit{First-order logic} extends propositional logic by introducing variables and quantifiers (like “for all” and “there exists”). This lets us make statements about objects and their properties, not just whole propositions.

    \item \textit{Second-order logic} goes even further, allowing quantification not only over objects but also over properties or sets of objects.
\end{itemize}



Propositional logic provides the foundation, but first-order and second-order logics are more expressive and powerful. For new students, propositional logic is the best place to start because its rules are simple, and it forms the basis for understanding more advanced topics. In later sections, I will discuss formal proofs and logical systems, and show how these principles connect to calculus and linear algebra. The aim is to provide a strong logical foundation to help you understand more complex mathematical ideas as you continue through this website.




\newpage

\section{Mathematical Notations and Terminology}
\subsection{What Are Propositions?}
In formal logic, a \textit{proposition} is a declarative statement that is either true or false, but not both. Propositions are the basic building blocks of logical reasoning and are used to construct more complex statements through logical connectives.

\textit{Propositional logic} focuses on the relationships between these propositions and how they can be combined using logical operators such as `and' (conjunction), `or' (disjunction), `not' (negation), `if...then' (implication), and `if and only if' (biconditional).



For example, take the following four different propositions represented by the variables \(r\), \(s\), \(c\), and \(w\):
\vspace{-0.5em}
\begin{align*}
    r\, &- \text{`It is raining'}\\
    s\, &- \text{`The sun is shining'}\\
    c\, &- \text{`It is cloudy'}\\
    w\, &- \text{`The grass is wet'}
\end{align*}

Here, each proposition gets a truth value. That is, \(r\), \(s\), \(c\), and \(w\) will be either `True' or `False'. We can build more complex propositions by joining these variables with logical operators like negation (\( \neg \)), conjunction (\(\land\)), disjunction (\(\lor\)), and implication (\(\rarrow\)) as shown below:
\vspace{-2.5em}
\begin{adjustwidth}{-6em}{0cm}
    \begin{align*}
        \neg r\, &- \text{It is not raining}\\
        r \land w\, &- \text{`It is raining and the grass is wet'}\\
        s \lor r\, &- \text{`The sun is shining or it is raining'}\\
        \neg (r \lor s) \rarrow c\, &- \text{`It is neither raining nor is the sun shining, which means it must be cloudy'}
    \end{align*}
\end{adjustwidth}

An easy way to remember the difference between \(\land\) for conjunction and \(\lor\) for disjunction is that \(\land\) looks like a capital `A' for `and'. 

\egblueboxs{PropEg}{
Consider the following propositions:
\begin{multicols}{2}
\begin{itemize}
    \item \(P(x): x\) is Prime
    \item \(O(x): x\) is odd.
\end{itemize}
\end{multicols}
We then have:
\begin{multicols}{2}
    \begin{itemize}
        \item \(P(11)\) is true
        \item \(\neg P(7)\) is false
        \item \(P(9) \land O(9)\) is false
        \item \(P(2) \lor O(2)\) is true
    \end{itemize}
\end{multicols}
}

\newpage

\subsection{Sets}
A set is a collection of objects considered as a single entity. The objects in a set are called its elements or members. Sets can contain any kind of object, such as numbers, symbols, or even other sets. One common way to formally define a set is to list its elements inside curly brackets. For example, the set \(\{1, 2, 3\}\) contains the elements \(1\), \(2\), and \(3\).

The order in which we list the elements does not matter, and repeating elements does not change the set. So, \(\{3, 1, 1, 1, 2\}\) is the same as \(\{1, 2, 3\}\). If we want to keep track of how many times each element appears, we use the term multiset instead of set. For instance, \(\{3\}\) and \(\{3, 3\}\) are different as multisets but are the same as sets. 

An infinite set has infinitely many elements. Since we cannot list all the elements of an infinite set, we often use the notation '\(\ldots\)' to indicate that the sequence continues forever.

We can use set-builder notation to describe a set containing elements according to some rule, which is often in the form \(\{n : [\text{rule about } n]\}\). For example, the set \(\{n \in \R : 2 \mid n\}\) represents the set of all real numbers such that \(n\) is divisible by \(2\).

\egblueboxs{SetsEg}{
Consider the following sets:
\begin{enumerate}
    \item The set of irrational numbers
    \item The set of negative integers greater than \(-6\)
    \item The set of integers that are multiples of \(4\) but not multiples of \(8\)
    \item The set of all numbers that can be written as \(3n - 2\) where \(n \in \Z\)
    \item The empty set
\end{enumerate}
Using set builder notation, we can define these sets as follows:
\begin{multicols}{2}
    \begin{enumerate}
        \item \(\{ x : (x \in \R) \land (x \notin \Q) \}\)
        \item \(\{ x \in \Z : -6 < x < 0 \}\)
        \item \(\{ x : (4 \mid x) \land (8 \nmid x) \}\)
        \item \(\{ 3n - 2 : n \in \Z \}\)
        \item Let \(x \in A\), and \(P(x)\) be false for all \(x\). Then, \(\emptyset = \{ x : P(x) \}\)
    \end{enumerate}
\end{multicols}
}


\newpage

\subsection{Set-theoretic Symbols}
In formal logic, set-theoretic symbols are used to represent sets and their relationships. If we take the sets, \(S = \{1, 2, 3\}\) and \(T = \{2, 3, 4\}\), we can use the following symbols to describe their relationships:

\paragraph{Is an Element Of.} The symbol \(\in\) denotes `is an element of'. If \(x \in A\), then \(x\) is an element of set \(A\). For example, \(2 \in S\) means that 2 is an element of set \(S\), and \(5 \notin T\) means that 5 is not an element of set \(T\). We use these symbols to express membership within sets.

\paragraph{Is a Subset Of.} The symbol \(\subseteq\) denotes `is a subset of'. If \(A \subseteq B\), then every element of set \(A\) is also an element of set \(B\). Note that \(A\) may be equal to \(B\). For example, if we have a set \(C = \{2, 3\}\), then \(C \subseteq S\) because all elements of \(C\) (which are 2 and 3) are also elements of \(S\). However, \(S \nsubseteq T\) because \(S\) contains the element 1, which is not in \(T\).

\paragraph{Is a Proper Subset Of.} The symbol \(\subset\) denotes `is a proper subset of'. If \(A \subset B\), then every element of set \(A\) is an element of set \(B\). However, \(A\) is not equal to \(B\). For example, using the same sets \(C\) and \(S\) as before, we can say that \(C \subset S\) because all elements of \(C\) are in \(S\), and \(C\) is not equal to \(S\).

It's important to note that the definition of this symbol (\(\subset\)) changes depending on the author. In some texts, \(\subset\) is used to denote `is a subset of', while in others, it denotes `is a proper subset of'. Throughout this document, the symbol \(\subseteq\) will be used to denote `is a subset of', and \(\subset\) will be used to denote `is a proper subset of'.

\paragraph{Union.} The symbol \(\cup\) denotes `union'. The union of sets \(A\) and \(b\) is the set of all elements that are in either \(A\) or \(B\). For example, the union of sets \(S\) and \(T\) is given by: \( S \cup T = \{ 1, 2, 3, 4 \} \).

\paragraph{Intersection.} 
The symbol \(\cap\) denotes `intersection'. The intersection of sets \(A\) and \(B\) is the set of all elements that are in both \(A\) and \(b\). For example, the intersection of sets \(S\) and \(T\) is given by: \(S \cap T = \{ 2, 3 \}\). If there are no common elements between the two sets, their intersection is the empty set, denoted by \(\emptyset\). For example, if we have sets \(U = \{ 5, 6 \}\) and \(V = \{7, 8\}\), then \(U \cap V = \emptyset\).

\paragraph{Set Difference.}
The symbol \(\setminus\) denotes `set difference'. The set difference of sets \(A\) and \(B\) is the set of all elements that are in \(A\) but not in \(B\). For example, the set difference of sets \(S\) and \(T\) is given by: \(S \setminus T = \{1\}\), since 1 is in \(S\) but not in \(T\). If we take the sets \(U = \{ 5, 6, 7, 8 \}\) and \(V = \{ 6, 7\}\), then the set difference is \(U \setminus V = \{ 5, 8 \}\).

\paragraph{Such That.}
The symbols \(\mid\) and \(:\) both denote `such that'. They are often used in set-builder notation to define a set by specifying a property that its members must satisfy. For example, the set of all even natural numbers can be defined as: \( \{ x \in \N_0 : x \text{ is even}\} \). This reads as `the set of all natural numbers \(x\) such that \(x\) is even'.

In this document, the symbol \(:\) will be primarily used to denote `such that'. This is to avoid confusion, as the vertical bar symbol \(\mid\) is also used in number theory to denote `divides'. For example, \(a \mid b\) means That \(a\) is a divisor of \(b\) (i.e., \(b\) is divisible by \(a\)).

\newpage

\subsection{The Set of Real Numbers (and its Subsets).}
The set of real numbers and its subsets are represented by the following:

\begin{itemize}
    \item \(\R\): The set of real numbers.
    \item \(\Q\): The set of rational numbers.
    \item \(\Z\): The set of integers.
    \item \(\N_0\): The set of natural numbers.
    \item \(\N\): The set of natural numbers excluding zero.
\end{itemize}

There is no specific symbol for the set of irrational numbers, but they can be represented as \(\R \setminus \Q\), which denotes the set of all real numbers that are not rational. Some authors may use custom definition, such as \(\mathbb{P} = \{ x : x \in \R \land x \notin \Q \}\), to denote the set of irrational numbers. However, this document will use the set difference notation \(\R \setminus \Q\) for the irrational numbers.

Note that some authors define the set of natural numbers \(\N\) to include 0, while others define it to start from 1. In this document, I'll use \(\N_0\) to denote the set of natural numbers including 0, and \(\N\) to denote the set of natural numbers excluding 0. When referring to the `natural numbers', it should be implied that I am referring to \(\N_0\) unless context implies otherwise. If there is any ambiguity, clarification will be provided.

The set of integers \(\Z\) includes the natural numbers (including 0), and the negative numbers. Using some basic notation, we can define the following variations of the set of integers:

\begin{itemize}
    \item \( \Z_{<0} = \{ x \in \Z : x < 0 \} \): The set of negative integers (i.e., \( \{-1, -2, -3, \ldots\} \)).
    \item \( \Z_{\leq 0} = \{ x \in \Z : x \leq 0 \} \): The set of negative integers including zero (i.e., \( \{0, -1, -2, -3, \ldots\} \)).
\end{itemize}

Similarly, we can define variations of the set of rational numbers:


\begin{itemize}
    \item \(\Q_{> 0} = \{ x \in \Q : x > 0 \}\): The set of positive rational numbers.
    \item \(\Q_{\geq 0} = \{ x \in \Q : x \geq 0 \}\): The set of non-negative rational numbers.
    \item \(\Q_{< 0} = \{ x \in \Q : x < 0 \}\): The set of negative rational numbers.
    \item \(\Q_{\leq 0} = \{ x \in \Q : x \leq 0 \}\): The set of non-positive rational numbers.
\end{itemize}

Similar notation can be applied to the set of real numbers \(\R\) as well:

\begin{itemize}
    \item \(\R_{> 0} = \{ x \in \R : x > 0 \} = (0, \infty)\): The set of positive real numbers.
    \item \(\R_{\geq 0} = \{ x \in \R : x \geq 0 \} = [0, \infty) \): The set of positive real numbers including zero.
    \item \(\R_{< 0} = \{ x \in \R : x < 0 \} = (-\infty, 0)\): The set of negative real numbers.
    \item \(\R_{\leq 0} = \{ x \in \R : x \leq 0 \} = (-\infty, 0]\): The set of negative real numbers including zero.
\end{itemize}

Some authors may use other notation notation to represent subsets (for example, \( \R^+ \) can represent the non-negative real numbers). However, I've decieded to use the above notation to ensure carity as to whether 0 is or is not included in the subsets. Set builder notation will be used if clarification is needed.

\newpage

\subsection{Universal and Existential Quantifiers}
Quantifiers are symbols that tell us how many members of a set or group a statement applies to. Instead of talking about just one particular object, quantifiers allow us to make statements about all objects in a group or about the existence of at least one object with a certain property.

\paragraph{Universal Quantifier.}
The universal quantifier is denoted by the symbol \(\forall\) and is read as `for all' or `for every'. It is used to indicate that a statement applies to all elements in a particular set or domain. For example, we'll take the following propositions:

\begin{itemize}
    \item \(P(x)\): \(x\) is prime greater than 2.
    \item \(O(x)\): \(x\) is odd.
\end{itemize}

These are both propositions because they can be either true or false depending on the value of \(x\). For example, if \(x = 5\), then \(P(5)\) is true and \(O(5)\) is true. However, if \(x = 4\), then \(P(4)\) is false and \(O(4)\) is false. Using the universal quantifier, we can express the statement that all prime numbers greater than 2 are odd: \(\forall x (P(x) \rarrow O(x))\), which reads: `for all values of \(x\), if \(x\) is prime greater than 2, then \(x\) is odd.'

\paragraph{Existential Quantifier.}
The existential quantifier is denoted by the symbol \(\exists\) and is read as `there exists' or `there is at least one'. It is used to indicate that there is at least one element in a particular set or domain that satisfies a given property. For example, the statement \(\exists x ( x \in \R \land x^2 = 4)\) means `there exist a number \(x\) where \(x\) is a real number and \(x\) squared equals 4', or more naturally, `there exists a real number whose square is 4'. In this case, the values \(x = 2\) and \(x = -2\) both satisfy the property.

\egblueboxs{QuantifersEg}{
    Consider the following sets, given the universal set is \(\R\):
    \begin{multicols}{2}
        \begin{enumerate}
            \item \(\Big\{ x : \forall n\Big((n \in \N) \rarrow (x = 3n - 2) \Big) \Big\}\)
            \item \(\Bigg\{ x : \exists n\Bigg( (n \in \Z) \land \Big((x = n^2) \land (x < 0)\Big)\Bigg) \Bigg\}\)
            \item \(\Big\{ x : \forall n\Big( (n \in \N \land x = n) \rarrow n^2 = x\Big) \Big\}\)
            \item \(\Big\{x : \exists! n \in \Z\Big((n < 8) \land (n^2 = x)\Big)\Big\}\)
        \end{enumerate}
    \end{multicols}

    We can simplify these sets by doing the following:
    \begin{enumerate}
        \item If the statement is true for \textit{all} natural numbers, then \(x\) must equal \(3n - 2\) for every \(n \in \N\). However, \(3n - 2\) takes on multiple values (\(1,~3,~7,~\cdots\)), and there is no value \(x\) that can equal them all at once. Thus, we get the empty set.
        \item The square of an integer is never negative. Thus, you get the empty set.
        \item If \(x \notin \Z\), then \(n \in \Z\) is false. So the implication is true for all  \(n \in \R \setminus \Z\), and thus every non-natural real \(x\) is included. If \(x \in \N\), then we need \(x = n \rarrow x = n^2\). This is only true for \(0\) and \(1\). Thus, we get: \( (\R \setminus \Z) \cup \{0, 1\}\).
        \item We need exactly one integer \(n\) such that both \(n < 8\) and \(n^2 = x\) hold. For \(x=0\), only \(n=0\) satisfies \(n^2=0\), so \(x=0\) is valid. For any positive square \(x = k^2\) with \(k < 8\), both \(n = k\) and \(n = -k\) satisfy \(n^2 = x\) and \(n < 8\), so there are two solutions, which is not allowed. If \(k > 8\), then \(+k\) does not satisfy \(n < 8\), but \(-k\) does, so there is exactly one solution.
        
        Therefore the set contains \(0\) and all squares \(k^2\) where \(k\in\Z\) and \(k \leq -8\). Because \(k^2 = (-k)^2\), an equivalent set is \(\{0\} \cup \{ n^2 : n \in \Z \land n \geq 8 \}\).
    \end{enumerate}
}

\section{Propositional Formulas}

\subsection{Writing Propositional Formulas}
\subsubsection{Formal Definition and Notation}

A propositional formula is an expression made up of propositions which are combined using logical connectives. This is simiar to how algebraic expressions are made up with terms (like \(2x\) or \(5\)), which are combined with operations (like taking \(2x + 5\)). Once we assign truth values to all the variables in the propositional formulas, the formula itself has a definite truth value: either true or false.

More formally, a propositional formula is not a proposition itself, but rather a formal expression that stands for a proposition. This is similar to how the algebraic expression represents a value, but is not the value itself.

When we specify values for all the variables in a propositional formula, we determine its unique truth value. Propositional formulas give us a systematic way to represent and analyse logical relationships between statements.

When we combine propositions with logical connectives, we must follow strict rules. These rules tell us exactly which combinations of symbols count as well-formed formulas. The symbols used to write propositional formulas are listed below.

\begin{itemize}
    \item Variables: \(p_1, \, p_2, \, p_3, \, \ldots\);
    \item Propositional connectives: \(\neg\), \(\land\), \(\lor\), \(\rarrow\), \(\ioi\); and
    \item Brackets: ( and ).
\end{itemize}

\textdef{Definition \defcs.}
A propositional formula is an expression made up of certain symbols, arranged according to specific rules. These rules tell us exactly which combinations of symbols count as well-formed formulas. The symbols used to write propositional formulas are listed below.

\vspace{-0.5em}

\begin{itemize}
    \item Any variable of the form \(p_i\), where \(i \geq 1\), is a propositional formula;
    \item If \(A\) is a propositional formula, then \(\neg A\) is also a propositional formula; and
    \item If \(A\) and \(B\) are propositional formulas, then the expressions \((A \lor B)\), \((A \land B)\), \((A \rarrow B)\) and \((A \sioi B)\) are also propositional formulas.
\end{itemize}

These are the only ways we are allowed to construct propositional formulas. Brackets must be included exactly as shown so that the meaning is clear. Any expression that does not follow these rules, such as leaving out or adding extra brackets, is not a well-formed formula and should not be used in formal logic.

\egblueboxss{Examples Of Propositional Formulas}{
    The main idea at this point is that a propositional formula is defined by how its symbols are arranged according to the rules above. At this point, the formula does not need to have any meaning assigned; it is enough that it is well-formed.

    \vspace{-1em}
    
    \begin{multicols}{2}
        \begin{itemize}
            \item \(p_i\)
            \item \(((\neg p_1 \land p_2) \ioi p_3)\)
            \item \( \neg (p_1 \rarrow \neg \neg p_2) \)
            \item \((p_1 \rarrow p_2) \ioi (p_2 \land p_3)\)
        \end{itemize}
    \end{multicols}

    \vspace{-1em}

    However, the expression \(p_1  \land p_2\) is not a well-formed propositional formula because it is missing the required brackets. The expression \((\neg p_3)\) is also not well-formed, since it has extra brackets that are not needed. Finally, \((p \lor r)\) is not a well-formed propositional formula because it uses variables that are not of the form \(p_i\).
}

\subsubsection{Informal Notation}

The formal definition uses a lot of brackets to make the structure of formulas clear, but all these brackets can make formulas harder to read. It is common to leave out some brackets when writing formulas informally, as long as the meaning is still clear. In particular, the outermost brackets are often left out. It is important to be able to put the brackets back in by following the rules of precedence, which I will describe below.

\begin{itemize}
    \item The negation sign has the highest precedence. Just as 'not' can change the meaning of a whole sentence, the negation sign changes the truth value of the entire expression it applies to before any other operation.
    \item The connectives \(\land\) and \(\lor\) have the second-highest precedence.
    \item The connectives \(\rarrow\) anad \(\ioi\) have the lowest precedence
    \item If we leave out brackets between connectives that have the same level of precedence, we group them from right to left. For example, the expression \( p \lor q \lor r\) is read as \( (p \lor (q \lor r))\).
\end{itemize}

\egblueboxss{Formal vs Informal Formulas}{
    The table below shows the difference between a formally-written formula, and its informal abbreviation.
    \begin{table}[H]
        \begin{tabular}{|c|c|}
            \hline \textbf{Informal Abbreviation} & \textbf{Actual Formula} \\\hline
            \(\neg p_1 \land p_2 \ioi p_3\) & \(((\neg p_1 \land p_2) \ioi p_3)\)\\\hline
            \(\neg p_1 \land p_2 \lor \neg p_3\) & \((\neg p_1 \land (p_2 \lor \neg p_3))\)\\\hline
            \(\neg p_1 \lor p_2 \lorarrow p_3 \sioi p_4\) & \(((\neg p_1 \lor p_2) \lorarrow (p_3 \sioi p_4))\) \\\hline
            \(p_1 \lorarrow p_2 \sioi p_3 \land p_4\) & \((p_1 \lorarrow (p_2 \sioi (p_3 \land p_4)))\) \\\hline
        \end{tabular}
    \end{table}
}

Besides omitting brackets, it is also common to use variable names like \(p, q, r, \ldots\) instead of \(p_1, p_2, p_3, \ldots\) to make formulas easier to read. Formally, only variables of the form \(p_i\) are allowed, but using single-letter names is a helpful informal convention.

The connectives \(\land\), \(\lor\) and \(\sioi\) are associative, so associating from right to left does not change their meaning. However, \(\rarrow\) is not associative. For example, \(p \rarrow (q \rarrow r)\) is not the same as \((p \rarrow q) \rarrow r\).

However, \(p \rarrow (q \rarrow r)\) is truth-functionally equivalent to \(p \land q \rarrow r\). This is a useful way to write formulas and is one reason for associating from right to left.

To avoid confusion, do not write expressions such as \( A \lor B \land C\) or \(A \lorarrow B \sioi C\) without brackets. According to the conventions here, these mean \((A \lor (B \land C))\) and \((A \lorarrow (B \sioi C))\), but other authors might use different conventions.





\newpage
\pagenumbering{gobble}\section*{} \ \vfill
\textbf{Document Management} \\
Version: 1.0 \\
Created: 19 February 2026 \\ 
Reviewed: 22 March 2026 \\
Contact: archimedes314notes@gmail.com


\end{document}
